\paragraph{Beyond an invariant channel.} Corollary~\ref{cor:recurrence} gives a closed scalar formula because all increments point in the same direction. General stable maps can rotate, cancel, or transiently amplify a readout. Their completion criterion is still explicit as a sequence of matrix partial sums, with a quantitative bound on the unresolved tail.

\begin{proposition}[General projected first passage and finite certification]\label{prop:generalpassage}
Under Theorem~\ref{thm:reciprocal}, put $p_0=L_{PQ}d$, let $v$ be a fixed readout, and set $q_A=\|A\|<1$. Define
\[
z_m=v^\top\sum_{k=0}^{m}A^kp_0,
\qquad z_\infty=v^\top(I-A)^{-1}p_0,
\qquad E_m=\frac{\|v\|\|p_0\|q_A^{m+1}}{1-q_A}.
\]
For a positive signed threshold $\theta$, the first attained event is $\tau_f+m_*T_{\rm loop}$, where $m_*=\min\{m\ge0:z_m\ge\theta\}$, or it never occurs if this set is empty. The tail satisfies $|z_\infty-z_m|\le E_m$. If $z_\infty>\theta$, completion is necessarily finite; any $m$ with $E_m<z_\infty-\theta$ is an upper bound on the crossing index. If all values through $M$ are below threshold and $z_\infty+E_M<\theta$, completion never occurs. Values $z_\infty\le\theta$ alone do not exclude an earlier transient crossing. For $n_P$ receiver coordinates, the readout is identically zero if and only if $v^\top A^kp_0=0$ for $k=0,\ldots,n_P-1$.
\end{proposition}
\begin{proof}
Equation~\eqref{eq:echo} is constant on each interval between consecutive return times, including its attained left endpoint. Applying $v^\top$ gives $z_m$ on the $m$th interval, proving the minimum-index rule without an invariance or positivity assumption. Absolute convergence yields
\[
|z_\infty-z_m|\le\|v\|\sum_{k=m+1}^{\infty}\|A\|^k\|p_0\|
=E_m.
\]
Since $E_m\to0$, if $z_\infty>\theta$ there is a finite $m$ with $z_m\ge z_\infty-E_m>\theta$. The first crossing is no later than this index. If $q_A=0$, the tail is zero and the first partial sum is already the limit. For every $m\ge M$, $E_m\le E_M$, so $z_m\le z_\infty+E_M<\theta$ certifies exclusion of all later crossings after the finite checked prefix. Equality in either bound does not provide the corresponding strict certification.

An identically zero readout has zero successive differences, so every $v^\top A^kp_0$ vanishes. Conversely, the Cayley--Hamilton identity expresses $A^{n_P}$ as a linear combination of lower powers. Multiplication by further powers and induction express all higher powers using $I,A,\ldots,A^{n_P-1}$. If the first $n_P$ scalar coefficients vanish, all subsequent ones vanish, and so does every partial sum. This proves the finite support test. None of these arguments requires $p_0$ to be an eigenvector.
\end{proof}

For example, take $A=\left(\begin{smallmatrix}0&0.5\\-0.5&0\end{smallmatrix}\right)$ and $p_0=v=e_1$. The first readout is 1, while its limit is 0.8; threshold 0.9 is crossed at direct arrival despite lying above the limit. Threshold 0.8, equal to the limit, is also crossed at direct arrival. Thus the scalar invariant-channel exclusion at equality is not general. With $p_0=e_1$, $v=e_2$, and $A=\left(\begin{smallmatrix}0&0\\0.5&0\end{smallmatrix}\right)$, the direct readout is zero and the first return is 0.5. Rotation into a readout can therefore create a later contribution even without scalar accumulation. For absolute-amplitude completion, replace $z_m\ge\theta$ by $|z_m|\ge\theta$; the same tail bound applies with $|z_\infty|$, while the signed example and minimum must be evaluated with the chosen rule.
